\scanpage{48}

Le groupe $\mathbb Z/2\mathbb Z$ opère sur $X'/X$, via

\[
\left\{\begin{aligned}
&x\longmapsto\deleted{\unsure{T-x}}\bar x=T-x\quad\text{sur }\underline P_1\deleted{\unsure{^{*}}}\qquad(\text{auto})\\
&x+\bar x=T\qquad\text{si }x\in\Gamma\underline P_1
\end{aligned}\right.
\]

ou encore

\[
\left\{\begin{aligned}
&x\longmapsto\bar x=\psi(x)T-x\quad\text{sur }E\\
&\text{i.e.}\quad x+\bar x=\psi(x)T
\end{aligned}\right.
\]

On a

\[
\left\{\begin{aligned}
&\left.\begin{aligned}\psi(T)&=T\\\psi(N)&=N\end{aligned}\right\}\quad\nwarrow\\
&Q(\deleted{\unsure{\psi(x)}}\bar x)=Q(x)
\end{aligned}\right.
\]

Sur $\mathscr A$, l'involution donnée par

\[
x\longmapsto\bar x=\operatorname{Tr}(x)-x\qquad\text{i.e.}\qquad x+\bar x=\operatorname{Tr}(x).1
\]

Considérons $X_0=V(2)$, \deleted{\unsure{sur} \unsure{$T_0$} $\unsure{T\vert V(2)|N(x)}\in\Gamma\underline L$ et $T_0,N_0,$} \inserted{et $E_0,L_0,T_0,N_0,\unsure{Q_0}$ \inserted{$\delta_0$}} les restrictions à $X_0$. On a \deleted{$\unsure{T_0\in\Gamma|T_0\in E_0}$} alors

\[
T_0\in\Gamma\underline L_0\quad(\Gamma E_0)
\]

\[
\text{et}\qquad T_0^2=\delta_0\quad(=-\delta_0)\ \text{sur }X_0
\]

et $=$ mieux~: si $T_0$ est relevé loc\textsuperscript{t}. \inserted{en section $\tau_1$ de $E$} sur $X_1=V(4)$, \deleted{\unsure{on aura}} (p. ex. $\tau_1=T|X_1\overset{\mathrm{déf}}=T_1$) on aura

\[
\tau_1^2=\delta_1\quad(\overset{\mathrm{déf}}=\delta|X_{\unsure{4|1}})
\]

(NB cela ne dépend pas du choix de ce relèvement, car $(\tau_1+2\alpha)^2\equiv\tau_1^2\ (4)$) On l'exprimera par l'écriture

\[
\boxed{T_0^2\equiv\deleted{\unsure{-}}\delta\quad(4)}
\]

Ceci posé, si $x$ section de $X'/X$, $z$ la section corr. de $\underline L$, via $z=2x-T$, \deleted{\unsure{$x_0$|$z_0$}}, $z_0$ les restr. à $X_0$, on aura

\[
\left\{\begin{aligned}z_0&=T_0\\z^2&=\delta\end{aligned}\right.
\]

i.e $z$ est une solution de $z^2=\delta$ \inserted{sur $X$} qui prolonge la solution déjà donnée \unsure{$z_0$|$T_0$} sur $X_0$.

% archive sheet number: 47
